\section{Conclusion}
\label{sec:conclusion}

We introduced the \emph{GLPE family} of intrinsic-reward methods---\textbf{GLPE} and \textbf{GLPE (no gate)}---for exploration and control with on-policy optimization. Both variants couple \emph{feature-space impact} with \emph{region-local learning progress} computed from an online partition of the learned latent space (Section~\ref{sec:method}). This combination yields an intrinsic signal that is (i) local, reflecting where the dynamics model is currently improving, (ii) scale-stable through component-wise normalization, and (iii) easy to integrate into PPO via a clipped shaping term and a simple annealing schedule.

Across the five-task benchmark suite and evaluation protocol in Section~\ref{sec:experimental_setup}, the results in Section~\ref{sec:results} support \textbf{GLPE (no gate)} as the most reliable default: it remains consistently competitive with the strongest intrinsic baselines across environments under both step-normalized and time-normalized summaries, while avoiding gating-specific thresholds and overhead. \textbf{GLPE} adds a region-specific gate that suppresses intrinsically rewarding but unproductive experience, and is best interpreted as an optional guardrail: it is most beneficial in sparse-reward regimes where intrinsic shaping can otherwise drift toward high-error, low-progress regions, and it typically behaves similarly to GLPE (no gate) when gating rarely activates. The intrinsic cosine taper further reinforces this interpretation by making shaping transient, supporting early exploration without dominating late-stage optimization.

The GLPE family also highlights practical tradeoffs. Gating introduces additional computation and thresholding choices, and its conservative suppression can be counterproductive in high-variance settings where progress is already seed-limited. More broadly, both GLPE variants inherit the usual sensitivities of representation-based intrinsic motivation, including dependence on the learned feature space and on the granularity of the online partition. Future work could focus on reducing gating overhead without altering its semantics, improving adaptive region construction, and extending the same impact-plus-progress principle to higher-dimensional observation modalities and alternative training regimes.
